
The results in this paper separate two things that are easy to
conflate when picking an accelerator: what the hardware is capable of,
and what a given program actually gets from it. On every measure we
took, the single v5e chip's raw capability was never in question. What
varied, sometimes by an order of magnitude, was how much of that
capability the default execution path exposed, and closing that gap
required knowing which JAX transformation to reach for
(Section~\ref{sec:parallelism}), understanding what a profiler trace
was and was not measuring (Section~\ref{sec:profiling}), and pricing a
slice by whether it was kept busy rather than by its hourly rate
(Section~\ref{sec:cost}). None of this required modifying AlphaFold2's
own source.

\subsection{Steady state versus cold start}
\label{sec:disc-when}

The single-chip comparison in Section~\ref{sec:results-hardware}
favors the TPU by a wide margin at steady state, but the same section
shows a raw 59.1$\times$ gap between the TPU's cold and steady-state
call, uncorrected because no TPU profiler-overhead log survives to
correct it. Whatever the true corrected value, the qualitative point
carries: a workload dominated by one-off calls at varying input shapes
pays the compilation cost on nearly every call, while a workload that
reuses one compiled shape repeatedly, such as sustained serving of a
fixed-length input, amortizes it away. This paper measures a
steady-state advantage, and nothing here tests whether it holds for a
workload that never reaches steady state.

\subsection{The cost of the default path}
\label{sec:disc-default}

A user who runs our unmodified single-query forward-pass driver on an
8-chip TPU slice gets the result Section~\ref{sec:parallelism}
describes: one chip working, seven idle, with no error or warning that
anything is wrong. We did not test every production entry point, but
nothing in AlphaFold2's default configuration asks for anything
different. Section~\ref{sec:cost}
shows what this costs in money as well as in throughput: the pod bills
for 87.5\% capacity that produced no output. No bug is involved:
on our reading of the evidence, this is the documented behavior of
\texttt{jax.jit} and of an automatic partitioner given a model with no
sharding annotations. The gap this paper
measures is therefore not primarily a hardware limitation or a flaw in
AlphaFold2's implementation, but a gap in what a user needs to know
about JAX's parallelism primitives before an 8-chip rental delivers
8 chips of work.
